\documentclass[12pt]{article}
\usepackage[utf8]{inputenc}
\usepackage[T1]{fontenc}
\usepackage[margin=1in]{geometry}
\usepackage{setspace}
\usepackage{xcolor}
\usepackage[colorlinks=true, urlcolor=blue!60!black, linkcolor=blue!60!black]{hyperref}
\setlength{\parindent}{0em}
\setlength{\parskip}{0.6em}
\renewcommand{\baselinestretch}{1.0}
\begin{document}
\begin{center}
{\LARGE \textit{Dobbs v. Jackson Women's Health Organization}, 597 U.S. 215 (2022)\par}
\vspace{0.5em}
{\large Case Brief --- Two-Page Edition\par}
\end{center}

\section{Memory Jogger}
\textit{Roe} and \textit{Casey} are overruled: the Constitution confers no right to abortion, and the issue returns to the people and their elected representatives (\href{https://www.courtlistener.com/opinion/6481357/dobbs-v-jackson-womens-health-organization/}{CourtListener}).

\section{Detailed Case Facts}
Mississippi's 2018 Gestational Age Act banned abortion after 15 weeks, with exceptions only for medical emergency or severe fetal abnormality. Jackson Women's Health Organization, the State's only abortion clinic, sued. Under \textit{Roe} and \textit{Casey}, States could not ban abortion before viability (about 23 to 24 weeks).

\section{Procedural History}
The district court enjoined the Act, and the Fifth Circuit affirmed. The Supreme Court granted certiorari on whether all pre-viability prohibitions are unconstitutional and reversed on June 24, 2022.

\section{Judicial Votes}
6--3 to uphold the ban, 5--4 to overrule \textit{Roe} and \textit{Casey}. Alito wrote for Thomas, Gorsuch, Kavanaugh, and Barrett; Thomas and Kavanaugh concurred; Roberts concurred in the judgment; Breyer, Sotomayor, and Kagan jointly dissented.

\section{Holding}
No constitutional right to abortion exists. Unenumerated rights must be ``deeply rooted in this Nation's history and tradition'' under \textit{Glucksberg}; abortion is not. Abortion laws receive rational-basis review, which the 15-week ban satisfies.

\section{Analysis of Opinions}
\textbf{Majority:} History from the common law through 1868, when three-quarters of the States banned abortion at all stages, shows no such right; abortion is ``critically different'' from other privacy precedents because it destroys ``fetal life.'' Five \textit{stare decisis} factors, including \textit{Roe}'s ``exceptionally weak'' reasoning and the absence of concrete reliance, supported overruling. \textbf{Thomas:} reconsider all substantive due process precedents. \textbf{Kavanaugh:} the Constitution is ``neither pro-life nor pro-choice.'' \textbf{Roberts:} uphold the ban without overruling \textit{Roe} ``all the way down to the studs.'' \textbf{Dissent:} the ``1868 test'' consigns women to ``second-class citizenship,'' and only the Court's composition had changed.

\section{Examples: Future Applications}
\textbf{Same-Side:} a total ban from fertilization survives rational-basis review. \textbf{Opposite-Side:} a felony for traveling out of state for an abortion likely violates the right to interstate travel. \textbf{Fence-Sitter:} a law treating IVF embryos as ``unborn human beings'' tests \textit{Dobbs}'s history method against procreative liberty.

\section{Critique}
Critics argue that tying liberty to 1868 entrenches women's historical exclusion and that the \textit{Glucksberg} method cannot coherently spare \textit{Griswold} or \textit{Obergefell}. Some originalists fault the Court for keeping substantive due process at all, while minimalists side with Roberts.

\section{Key Quotations}
\begin{quote}
``We hold that \textit{Roe} and \textit{Casey} must be overruled. The Constitution makes no reference to abortion, and no such right is implicitly protected by any constitutional provision, including the one on which the defenders of \textit{Roe} and \textit{Casey} now chiefly rely---the Due Process Clause of the Fourteenth Amendment.''
\end{quote}
\noindent --- Justice Alito, for the Court, 597 U.S. 215.

\end{document}
